\scanpage{50}
% loc: continuation of item 5
Par contre, sans hypothèse sur $k$ sauf d'être \uline{alg. clos}, il semble encore raisonnable d'espérer que (4.2) soit tjrs pleinement fidèle.
% loc: left oval
% slanted
\margin{réseau $+$ connexion}
On est ramené pour ceci au cas $k=\mathbb C$\unsure{..|.} C'est une conséquence de la conjecture ``du double réseau'' de 4) $+$ Tate ou Hodge, \unsure{préc.|précisé} \unsure{à|:} 6)

% loc: item 6, main line and interlinear insertion
6) Si $S$ de type fini \inserted{sur \deleted{$\unsure{\mathbb Q}$} un corps de nbs}, on a le foncteur de 3), mais sans filtration sur $T_{\mathrm{DR}}(M)$. Il résultera\deleted{\unsure{it|t}}it de la conj. du double réseau \deleted{(qui, elle, \unsure{génér} \inserted{(de \unsure{3|?})}} que ce foncteur est pl. fidèle
% loc: closing NB
[\uline{NB} si $S$ est le spectre d'un corps de nb, cela se réduit \unsure{:|à} l'énoncé du double réseau\unsure{....})
% loc: paper watermark, read sideways and compared across sheets
% printed form
% paper watermark, visible letters in the direction of the paper lettering
\par\noindent\textsc{\unsure{trong|rong}}\hfill\textsc{Extra}\\
\unsure{ry|y} 1808\hfill\unsure{Mauno|Maun}
% archive sheet number: 49

